% §7 Open Problems
Our formal analysis reveals several open research directions that go beyond the scope of this paper.

\paragraph{Learned barge-in prediction.}
Current barge-in systems react to voice activity after it begins.
A predictive system would detect \emph{intent to speak} from prosodic cues (rising pitch, trailing cadence, filled pause) before the user actually starts speaking.
The agent can then pre-emptively silence its response.
This would effectively reduce the perceptual silence gap for interruption turns below zero---the agent stops before the user begins.
Formalising the prediction problem and establishing the tradeoff between false-positive barge-in rate and turn-taking smoothness is an open problem.

\paragraph{Safety in cascade voice systems.}
The cascade architecture introduces safety risks at each inter-component boundary.
ASR transcription errors can transform benign speech into harmful text inputs.
TTS injection attacks could exploit the text interface of the TTS component to produce audio content not intended by the LM.
Voice-channel jailbreaks may exploit prosodic or phonetic properties of the audio input that are discarded by ASR transcription.
A formal threat model for cascade voice systems, analogous to the text-based guardrail frameworks of \citet{dong2024building}, does not yet exist.

\paragraph{Voice-aware agent planning.}
The ReAct agent in our model is latency-unaware.
It selects tool calls and reasoning depth based solely on task requirements.
A voice-aware agent would reason about its remaining silence budget ($\tau_{human} - \Lambda_{\text{elapsed}}$) as a first-class planning constraint.
It would prefer shallower reasoning chains and pre-cached tool results when the budget is nearly exhausted.
Formalising this as a constrained planning problem and studying the quality-latency tradeoff it induces is an open direction.

\paragraph{Principled hybrid routing policies.}
Hybrid routing (Section~\ref{sec:design_levers}) requires a policy for classifying incoming turns as simple (cascade-suitable) or complex (native-suitable).
No principled approach to learning this routing policy currently exists.
Key challenges include defining the optimal routing objective (minimise cost subject to CCC satisfaction, or maximise quality subject to budget).
Handling turns that are initially classified as simple but escalate (the LM discovers mid-reasoning that a tool call is needed) is a further challenge.
Maintaining conversational context across pipeline switches (the native model and the cascade LM have different context representations) remains open.

\paragraph{Commit latency--accuracy Pareto frontier for streaming ASR.}
Theorem~\ref{thm:streaming_floor} shows that the streaming cascade floor $\Lambda^*_{SC}$
is dominated by the ASR commit latency $\delta_N$.
Reducing $\delta_N$ below 100~ms---the threshold at which the CCC becomes satisfiable
(Corollary~\ref{cor:streaming_ccc})---requires increasing the lookahead depth of the ASR
model, which decreases word error rate at the cost of increased $\delta_N$, or decreasing
lookahead, which reduces $\delta_N$ at the cost of higher WER.
No published work characterises this Pareto frontier for modern neural streaming ASR systems
paired with LLM-scale language models.
Formalising the tradeoff as a constrained optimisation---minimise expected $\delta_N$ subject
to a WER budget---and studying the optimal operating point across utterance types and domains
is an open problem with direct engineering impact.

\paragraph{Speculative cascade: early generation on partial transcripts.}
A natural extension of streaming ASR is to begin LM \emph{generation} (not just prefill)
on the stable ASR prefix before EOU is confirmed---analogously to speculative decoding
\citep{leviathan2023fast} but at the ASR hypothesis level.
Let $p_{accept}$ be the probability that the ASR does not revise any token that changes the
semantics of the LM's first response token.
The expected silence gap is:
\begin{equation}
  \mathbb{E}[\Lambda_{spec}] = p_{accept} \cdot \Lambda_{early}
    + (1 - p_{accept}) \cdot (\Lambda^*_{SC} + \lambda_{replan})
\end{equation}
Deriving $p_{accept}$ as a function of ASR model type and query distribution,
bounding the replan cost $\lambda_{replan}$,
and characterising the quality degradation from incorrect speculative generation
(especially when TTS has already begun audio playback from the wrong draft) are open problems.
The SDS literature \citep{schlangen2011general} treats revocability informally;
a formal quantitative treatment for transformer LM--streaming ASR pairs does not yet exist.

\paragraph{Probabilistic EOU detection as a design lever.}
Current EOU detection is reactive: the VAD window $d_{eou}$ is a fixed silence threshold.
Prosodic and syntactic predictors can estimate $P(\text{EOU} \mid \text{acoustic features},
C_{pre})$ before acoustic silence begins, triggering LM generation when this probability
exceeds a threshold $\theta$.
This converts $d_{eou}$ from a deterministic floor to a stochastic variable:
$d_{eou}(\theta) = \mathbb{E}[\text{anticipation horizon}] - P_{fal}(\theta) \cdot \lambda_{restart}$,
where $P_{fal}(\theta)$ is the false-EOU rate and $\lambda_{restart}$ is the cost of
restarting after a false trigger.
Characterising the optimal $\theta$ that minimises expected $\Lambda^*_{SC}$ subject to
a tolerable false-EOU rate, and establishing a revised impossibility condition when
$d_{eou}$ is treated as a stochastic design variable, is an open formal problem.
